In order to better aggregate results on NLP tasks and compare to popular models such as BERT and RoBERTa in a more systematic way, we also evaluate GPT-3 on a standardized collection of datasets, the SuperGLUE benchmark \cite{wang2019superglue} \citep{wang2019superglue} \citep{clark2019boolq} \citep{demarneffe:cb} \citep{roemmele2011choice} \citep{khashabi2018looking} \citep{zhang2018record} \citep{dagan2006pascal} \citep{bar2006second} \citep{giampiccolo2007third} \citep{bentivogli2009fifth} \citep{pilehvar2018wic} \citep{poliak2018dnc}. GPT-3’s test-set performance on the SuperGLUE dataset is shown in Table \ref{table:superglue}.  In the few-shot setting, we used 32 examples for all tasks, sampled randomly from the training set. For all tasks except WSC and MultiRC, we sampled a new set of examples to use in the context for each problem. For WSC and MultiRC, we used the same set of randomly drawn examples from the training set as context for all of the problems we evaluated.

We observe a wide range in GPT-3’s performance across tasks.  On COPA and ReCoRD GPT-3 achieves near-SOTA performance in the one-shot and few-shot settings, with COPA falling only a couple points short and achieving second place on the leaderboard, where first place is held by a fine-tuned 11 billion parameter model (T5). On WSC, performance is still relatively strong, achieving 80.1\% in the few-shot setting (note that GPT-3 achieves 88.6\% on the original Winograd dataset as described in Section \ref{section:Winograd-Style_Tasks}).  On BoolQ, MultiRC, and RTE, performance is reasonable, roughly matching that of a fine-tuned BERT-Large. On CB, we see signs of life at 75.6\% in the few-shot setting.

WiC is a notable weak spot with few-shot performance at 49.4\% (at random chance).  We tried a number of different phrasings and formulations for WiC (which involves determining if a word is being used with the same meaning in two sentences), none of which was able to achieve strong performance.  This hints at a phenomenon that will become clearer in the next section (which discusses the ANLI benchmark) -- GPT-3 appears to be weak in the few-shot or one-shot setting at some tasks that involve comparing two sentences or snippets, for example whether a word is used the same way in two sentences (WiC), whether one sentence is a paraphrase of another, or whether one sentence implies another. This could also explain the comparatively low scores for RTE and CB, which also follow this format. Despite these weaknesses, GPT-3 still outperforms a fine-tuned BERT-large on four of eight tasks and on two tasks GPT-3 is close to the state-of-the-art held by a fine-tuned 11 billion parameter model.

Finally, we note that the few-shot SuperGLUE score steadily improves with both model size and with number of examples in the context showing increasing benefits from in-context learning (Figure \ref{graph:superglue_analysis}). We scale $K$ up to 32 examples per task, after which point additional examples will not reliably fit into our context. When sweeping over values of $K$, we find that GPT-3 requires less than eight total examples per task to outperform a fine-tuned BERT-Large on overall SuperGLUE score.
